\documentclass[border=10pt]{standalone}
\usepackage{luatexja-fontspec}
\usepackage{tikz}
\usetikzlibrary{positioning, arrows.meta, calc}

% ヒラギノ丸ゴシックを指定
\setmainjfont{HiraMaruProN-W4}

\begin{document}
\begin{tikzpicture}[
    thick,                  % 全体の線を太めにする
    >=Stealth,              % 綺麗な矢じりを使用
    font=\large,            % フォントサイズを一括で大きくする
    % 標本化を表す白抜きの赤丸スタイル
    sample_dot/.style={draw=red!80!black, circle, minimum size=5pt, inner sep=0pt, fill=white, line width=1.2pt},
    % 追加された青い注釈用のスタイル
    text_blue/.style={text=cyan!80!black, align=center},
    line_blue/.style={draw=cyan!80!black, line width=1.0pt}
]

    % --- 1. 軸の描画 ---
    \draw [->] (-0.5, 0) -- (11.0, 0) node[right] {$t$};

    % --- 2. 連続信号の描画（黒い滑らかな曲線） ---
    \draw [black, line width=1.0pt, smooth, tension=0.6] plot coordinates {
        (-0.5, -0.4)
        (0.5, -0.9) (1.0, 0.4) (1.5, -0.3) (2.0, -0.2) (2.5, -0.8) (3.0, -0.4) 
        (3.5, 0.4) (4.0, 1.6) (4.5, 2.5) (5.0, 2.3) (5.5, 1.3) (6.0, 0.0) 
        (6.5, -0.6) (7.0, -1.0) (7.5, -0.5) (8.0, 0.4) (8.5, 0.1) (9.0, -0.4) 
        (9.5, 0.5) (10.0, -0.2)
        (10.5, -0.1)
    };

    % --- 3. 標本化パルス（赤い縦棒と丸）の描画 ---
    \foreach \x/\y in {
        0.5/-0.9, 1.0/0.4, 1.5/-0.3, 2.0/-0.2, 2.5/-0.8, 3.0/-0.4, 
        3.5/0.4, 4.0/1.6, 4.5/2.5, 5.0/2.3, 5.5/1.3, 6.0/0.0, 
        6.5/-0.6, 7.0/-1.0, 7.5/-0.5, 8.0/0.4, 8.5/0.1, 9.0/-0.4, 
        9.5/0.5, 10.0/-0.2
    } {
        \draw [red!80!black, line width=1.0pt] (\x, 0) -- (\x, \y);
        \node [sample_dot] at (\x, \y) {};
    }

    % --- 4. 凡例（右上のテキスト） ---
    \node [sample_dot] (legend_dot) at (8.0, 2.3) {};
    \node [red!80!black, right=0.15cm of legend_dot] {標本化};


    % --- 5. 手書きの青い注釈部分 ---
    % 【調整】重なりを回避するため、テキスト位置を (3.2, 3.2) から (2.3, 3.4) へ移動し、左上に逃がしました
    \node [text_blue, font=\normalsize] (info_text) at (2.3, 3.4) {最大の周波数が \\ $B$ [Hz] 未満};
    
    % 【調整】出発角度(out)と進入角度(in)を最適化し、黒い曲線の上空を綺麗にまたぐ軌道に変更しました
    \draw [line_blue, ->] (info_text.south east) to[out=-25, in=125] (4.0, 1.6);

    % サンプリング間隔を示す補助線
    \draw [line_blue, |-|] (4.0, -0.3) -- (4.5, -0.3);
    
    % 間隔の数式ラベル「1 / 2B [s]」
    \node [text_blue, below=0.1cm of {4.25, -0.3}, font=\normalsize] {$\displaystyle\frac{1}{2B}$ [s]};

\end{tikzpicture}
\end{document}

\documentclass[border=10pt]{standalone}
\usepackage{luatexja-fontspec}
\usepackage{tikz}
\usetikzlibrary{positioning, arrows.meta, calc}

% ヒラギノ丸ゴシックを指定
\setmainjfont{HiraMaruProN-W4}

\begin{document}
\begin{tikzpicture}[
    thick,                  % 全体の線を太めにする
    >=Stealth,              % 綺麗な矢じりを使用
    font=\large,            % フォントサイズを一括で大きくする
    % 標本化を表す白抜きの赤丸スタイル
    sample_dot/.style={draw=red!80!black, circle, minimum size=5pt, inner sep=0pt, fill=white, line width=1.2pt},
    % 追加された青い注釈用のスタイル
    text_blue/.style={text=cyan!80!black, align=center},
    line_blue/.style={draw=cyan!80!black, line width=1.0pt}
]

    % --- 1. 軸の描画 ---
    \draw [->] (-0.5, 0) -- (11.0, 0) node[right] {$t$};

    % --- 2. 連続信号の描画（黒い滑らかな曲線） ---
    \draw [black, line width=1.0pt, smooth, tension=0.6] plot coordinates {
        (-0.5, -0.4)
        (0.5, -0.9) (1.0, 0.4) (1.5, -0.3) (2.0, -0.2) (2.5, -0.8) (3.0, -0.4) 
        (3.5, 0.4) (4.0, 1.6) (4.5, 2.5) (5.0, 2.3) (5.5, 1.3) (6.0, 0.0) 
        (6.5, -0.6) (7.0, -1.0) (7.5, -0.5) (8.0, 0.4) (8.5, 0.1) (9.0, -0.4) 
        (9.5, 0.5) (10.0, -0.2)
        (10.5, -0.1)
    };

    % --- 3. 標本化パルス（赤い縦棒と丸）の描画 ---
    \foreach \x/\y in {
        0.5/-0.9, 1.0/0.4, 1.5/-0.3, 2.0/-0.2, 2.5/-0.8, 3.0/-0.4, 
        3.5/0.4, 4.0/1.6, 4.5/2.5, 5.0/2.3, 5.5/1.3, 6.0/0.0, 
        6.5/-0.6, 7.0/-1.0, 7.5/-0.5, 8.0/0.4, 8.5/0.1, 9.0/-0.4, 
        9.5/0.5, 10.0/-0.2
    } {
        \draw [red!80!black, line width=1.0pt] (\x, 0) -- (\x, \y);
        \node [sample_dot] at (\x, \y) {};
    }

    % --- 4. 凡例（右上のテキスト） ---
    \node [sample_dot] (legend_dot) at (8.0, 2.3) {};
    \node [red!80!black, right=0.15cm of legend_dot] {標本化};


    % --- 5. 【追加】手書きの青い注釈部分 ---
    
    % 「最大の周波数が B [Hz] 未満」のテキスト
    \node [text_blue, font=\normalsize] (info_text) at (3.2, 3.2) {最大の周波数が \\ $B$ [Hz] 未満};
    
    % テキストからサンプリング点（x=4.0の赤丸）へ向かう青い矢印
    \draw [line_blue, ->] (info_text.south east) to[out=-20, in=140] (4.0, 1.6);

    % サンプリング間隔を示す補助線（x=4.0 から x=4.5 の間）
    % |-| スタイルを使うことで、手書きの「H」のような両端の区切り線を表現できます
    \draw [line_blue, |-|] (4.0, -0.3) -- (4.5, -0.3);
    
    % 間隔の数式ラベル「1 / 2B [s]」
    \node [text_blue, below=0.1cm of {4.25, -0.3}, font=\normalsize] {$\displaystyle\frac{1}{2B}$ [s]};

\end{tikzpicture}
\end{document}

